\documentclass[11pt]{article}
\usepackage[letterpaper,margin=1in]{geometry}
\usepackage[T1]{fontenc}
\usepackage{lmodern}
\usepackage{amsmath,amssymb,amsthm}
\usepackage{booktabs,tabularx,array}
\usepackage[colorlinks=true,linkcolor=black,citecolor=black,urlcolor=blue]{hyperref}
\hypersetup{
  pdftitle={A nonuniform witness for Jensen--Toft 15.6},
  pdfauthor={Alex Chan}}

\newtheorem{theorem}{Theorem}
\newtheorem{lemma}{Lemma}

\title{A nonuniform witness for Jensen--Toft 15.6}
\author{Alex Chan\\[2pt]
  \small Harvard University and National Bureau of Economic Research\\[-1pt]
  \small \texttt{achan@hbs.edu}}
\date{October 2026}

\begin{document}
\maketitle

\begin{abstract}
We construct an eleven-vertex hypergraph with twenty edges of size four and
four edges of size five. It is three-chromatic, and its weight
\(\sum_{e}2^{-|e|}=11/8\) is below \(23/16\), the minimum weight of a
four-uniform hypergraph without Property B. Thus the uniform-attainment
question in Jensen and Toft's Problem 15.6 has a negative answer.
\end{abstract}

\section{The question and the construction}
A proper coloring of a hypergraph has no monochromatic edge. For \(c>1\), put
\[
  w_c(G)=\sum_{e\in E(G)}c^{-|e|},\qquad
  f(n)=\inf\{w_2(G):\chi(G)=3,\ |e|\ge n\text{ for all }e\in E(G)\}.
\]
Let \(m(n)\) be the smallest number of edges in an \(n\)-uniform hypergraph
without Property B. Jensen and Toft \cite[Problem~15.6, pp.~235--236]{JT}
ask whether \(f(n)\) is attained by an \(n\)-uniform three-chromatic
hypergraph. Since \(m(4)=23\) \cite{Ostergard}, every four-uniform
three-chromatic hypergraph has weight at least \(23/16\). Cherkashin
\cite[\S2]{Cherkashin} gave a nonuniform example of weight \(95/64\), which
remains above this benchmark.

\begin{theorem}\label{thm:main}
There is a three-chromatic hypergraph \(H\) with twenty four-edges and four
five-edges such that
\[
  f(4)\le w_2(H)=\frac{11}{8}<\frac{23}{16}=\frac{m(4)}{2^4}.
\]
Consequently no four-uniform hypergraph attains \(f(4)\).
\end{theorem}

\begin{proof}
Take four disjoint pairs
\[
 A=\{a_0,a_1\},\quad B=\{b_0,b_1\},\quad
 C=\{c_0,c_1\},\quad D=\{d_0,d_1\},
\]
and a disjoint three-element set \(U\). Write
\(T_{ijk}=\{a_i,b_j,c_k\}\) for \((i,j,k)\in\{0,1\}^3\). The edges of
\(H\) are exactly
\begin{align}
 P\cup S
  &\quad(P\in\{A,B,C,D\},\ S\in\tbinom{U}{2}),
       \tag{I}\label{eq:pairs}\\
 T_{ijk}\cup\{d_\delta\}
  &\quad(i+j+k\equiv0\pmod2,\ \delta\in\{0,1\}),
       \tag{II}\label{eq:even}\\
 T_{ijk}\cup D
  &\quad(i+j+k\equiv1\pmod2).
       \tag{III}\label{eq:odd}
\end{align}
These families have twelve, eight, and four edges, respectively, and are
disjoint. Thus \(H\) has eleven vertices, twenty four-edges, and four
five-edges.

Suppose that a proper red--blue coloring exists. Interchange the colors if
necessary so that at least two vertices of \(U\) are red. By
\eqref{eq:pairs}, each of \(A,B,C,D\) contains a blue vertex. If one of
\(A,B,C\) were wholly blue, choose blue vertices in the other two pairs and
then choose from the wholly blue pair the vertex that makes the index sum
even. A blue vertex of \(D\) would complete a blue edge of
\eqref{eq:even}. Hence each of \(A,B,C\) has one blue and one red vertex.

The blue transversal of \(A,B,C\) must have odd index sum, or a blue vertex
of \(D\) would complete an edge of \eqref{eq:even}. Its complementary red
transversal has even index sum, because all three indices change. A red
vertex of \(D\) would therefore complete a red edge of \eqref{eq:even}.
Both vertices of \(D\) are consequently blue. The odd blue transversal
together with \(D\) is now a blue edge of \eqref{eq:odd}, a contradiction.

For a proper three-coloring, color \(A,C,D\) with color 0, \(B\) with color
1, and \(U\) with color 2. Each edge of \eqref{eq:pairs} meets \(U\) and a
pair; each edge of \eqref{eq:even} or \eqref{eq:odd} meets both \(A\) and
\(B\). Thus \(\chi(H)=3\). Finally,
\[
 w_2(H)=20\cdot2^{-4}+4\cdot2^{-5}
       =\frac{11}{8}=\frac{22}{16}<\frac{23}{16}.
\]
The four-uniform lower bound stated above proves the claim.
\end{proof}

\begin{theorem}\label{thm:c}
For every \(4/3<c<2\), the infimum of the \(c\)-weights among
three-chromatic hypergraphs with edge sizes at least four is strictly less
than the weight of every four-uniform three-chromatic hypergraph.
\end{theorem}
\begin{proof}
The same \(H\) has
\[
 w_c(H)=20c^{-4}+4c^{-5}<23c^{-4}
 \quad\Longleftrightarrow\quad c>\frac43.
\]
Every four-uniform three-chromatic hypergraph has at least 23 edges.
\end{proof}

The construction provides an upper bound for \(f(4)\); it does not determine
its exact value. The finite encoding and verification files are prepared for
the companion repository at
\url{https://github.com/alexgocardinal/jt156-uniform-attainment}.
The repository link becomes active when the files are posted there.

\begin{thebibliography}{9}
\bibitem{JT}
T.~R. Jensen and B. Toft,
\emph{Graph Coloring Problems}, Wiley, New York, 1995, Problem~15.6,
pp.~235--236.

\bibitem{Ostergard}
P.~R.~J. \"Osterg\r{a}rd,
On the minimum size of 4-uniform hypergraphs without property B,
\emph{Discrete Appl. Math.} \textbf{163} (2014), no.~2, 199--204.
\url{https://doi.org/10.1016/j.dam.2011.11.035}.

\bibitem{Cherkashin}
D. Cherkashin,
On small non-uniform hypergraphs without property B,
\emph{Mathematics and Education in Mathematics} \textbf{52}
(2023), 71--75.
\url{https://doi.org/10.55630/mem.2023.52.71-75}.
\end{thebibliography}

\end{document}
